% introduction.tex -- Sequencer Paper (BNR-2026-009)

\section{Introduction}

The sequencer is the critical component of any Layer~2 rollup architecture: it receives
user transactions, orders them into blocks, and submits batches to the Layer~1 for
finality. In a zero-knowledge rollup, the sequencer's ordering decisions directly
determine the state transitions that the ZK prover must certify. The design of the
sequencer therefore affects not only throughput and latency but also the correctness
guarantees of the entire system.

A fundamental tension exists in sequencer design. Centralized (single-operator) sequencers
achieve the lowest latency and highest throughput because a single entity controls
transaction ordering without requiring consensus among multiple parties. However, this
centralization introduces a trust assumption: the operator could censor transactions,
reorder them for extractable value (MEV), or halt block production entirely. Decentralized
sequencer designs~\cite{taiko2024based} eliminate this trust assumption but inherit L1
block times and consensus overhead, reducing throughput by an order of magnitude.

The established solution across production rollups is to combine centralized sequencing
with a forced inclusion mechanism on L1~\cite{arbitrum2024architecture,
optimism2024derivation}. The sequencer operates efficiently under normal conditions,
while the L1 contract guarantees that any user can force-include a transaction if the
sequencer censors it, with a bounded deadline (typically 12--24 hours).

Basis Network is an enterprise blockchain platform deployed on Avalanche, targeting Latin
American industrial and commercial clients. The enterprise zkEVM L2 architecture
(established in BNR-2026-008~\cite{bnr2026008evm}) operates a per-enterprise chain model
where each client runs an independent L2 instance. This enterprise context introduces
unique constraints:

\begin{enumerate}
  \item \textbf{Zero-fee transactions}: The L1 operates with near-zero gas costs
    (1~wei base fee), eliminating fee-based transaction priority and rendering
    MEV extraction economically irrelevant.
  \item \textbf{Single known operator}: Each enterprise chain is operated by the
    enterprise itself, making the sequencer a trusted-but-verified component
    rather than an adversarial one.
  \item \textbf{FIFO fairness}: Without fee markets, the natural ordering policy
    is strict FIFO (first-in, first-out) by arrival time, providing deterministic
    ordering without priority auctions.
  \item \textbf{Sub-second block times}: Enterprise workloads (industrial
    traceability, commercial ERP) require low-latency confirmation, targeting
    1-second L2 block times.
\end{enumerate}

The research question addressed by this paper is:

\medskip
\noindent\textit{Can a single-operator sequencer produce L2 blocks every 1--2 seconds
with FIFO ordering while a forced inclusion mechanism via L1 guarantees censorship
resistance with maximum 24-hour latency?}
\medskip

We investigate this question through a comparative analysis of eight production L2
sequencer architectures, the design and implementation of a Go prototype
(${\sim}$500~LOC) featuring FIFO mempool, Arbitrum-style forced inclusion queue, and
tick-based block production, and comprehensive benchmarks across scaling, concurrency,
and adversarial scenarios.

The remainder of this paper is organized as follows. Section~II surveys related work
on L2 sequencer architectures and forced inclusion mechanisms. Section~III presents the
system model and design decisions. Section~IV describes the experimental methodology.
Section~V presents benchmark results across block production scaling, forced inclusion,
and concurrent access scenarios. Section~VI discusses implications, invariants, MEV
analysis, and limitations. Section~VII concludes with recommendations.
